\subsection{As duas aberturas de uma avaliação}
\label{sec:gbp-duas-aberturas}
\index{generalized bulletproofs@Generalized Bulletproofs!avaliacao mascarada@avaliação mascarada}

Vale a pena executar também as operações com as máscaras à vista.
Até aqui acompanhámos o número \(89\) desde os dois quadrados até
à última dobra. Em que pontos entra ele na prova? Surge uma vez como
valor de um compromisso sobre \(g,h\) e outra como coeficiente de
\(U\) num ponto formado com os geradores vectoriais. As duas
ocorrências têm de corresponder ao mesmo escalar que o provador
anuncia. Fazer passar a mesma avaliação de uma família de bases a
outra é justamente o que permite ao verificador exigir ambas as
coisas ao mesmo tempo.

Conservemos todos os valores anteriores em \(\mathbb F_{101}\).
Para acompanhar a primeira abertura, sejam \(7\) a máscara de
\(V\) e \(i+2\) a máscara do coeficiente de grau \(i\) para
cada \(i\in\mathcal I\). Estas escolhas servem para que o leitor
possa seguir as parcelas. Numa prova, as máscaras novas são
aleatórias. O verificador tem em mãos \(V\), com o qual obtém
\[
 B_4=72g+32V=(72+32\cdot25)g+(32\cdot7)h=64g+22h.
\]
Para o provador, \(64,22\) constituem a abertura desta combinação.
Ao verificador pertencem os pontos e os multiplicadores públicos
\(72,32\). A mesma operação de grupo dará \(B_4\) aos dois, ainda
que só um conheça os números pelos quais acabámos de a decompor.

Avaliemos no desafio \(\xi=3\). Podemos dispor os pontos a combinar
por ordem dos graus, juntando à direita as duas partes da sua abertura.
Cada coluna de contribuições está já reduzida módulo \(101\).
Designamos por \(\tau_4\) a máscara derivada, para escrever a conta
com a mesma notação em todos os graus.
\begin{center}
\begin{tabular}{c|c|rr|r|rr}
 Grau & Ponto & \(t_i\)&\(\tau_i\)&\(3^i\)&\(3^it_i\)&\(3^i\tau_i\)\\ \hline
 2&\(T_2\)&72&4&9&42&36\\
 3&\(T_3\)&85&5&27&73&34\\
 4&\(B_4\)&64&22&81&33&65\\
 5&\(T_5\)&35&7&41&21&85\\
 6&\(T_6\)&45&8&22&81&75\\
 7&\(T_7\)&82&9&66&59&89\\
 8&\(T_8\)&17&10&97&33&61\\
 9&\(T_9\)&59&11&89&100&70\\
10&\(T_{10}\)&7&12&65&51&73\\ \hline
 &\multicolumn{4}{r|}{Soma no corpo}&89&83
\end{tabular}
\end{center}
Por exemplo, na quarta potência temos \(81\cdot64=5184\), que
deixa resto \(33\), e \(81\cdot22=1782\), que deixa resto \(65\).
Os totais inteiros das duas colunas são \(493\) e \(588\), com
restos \(89\) e \(83\). Assim, o provador entrega \(\widehat t=89\)
e \(\tau_\xi=83\). O verificador multiplica cada ponto da segunda
coluna pela potência correspondente, soma e compara o resultado com
\(89g+83h\). Quem só vê a prova tem tudo quanto é necessário para
esta igualdade. Os valores e as máscaras de cada linha ficam com
o provador.

Passemos à outra abertura. O verificador parte dos pontos que
recebeu antes de conhecer \(\xi\). Tomemos as máscaras \(11,4,6\)
para \(A_I,A_O,S\), respectivamente, e \(13\) para \(K\).
Os pesos destes quatro pontos em \(P\) são as potências de \(3\)
que o quadro dos graus prescreveu. A soma das máscaras vale
\[
 \mu=9\cdot11+81\cdot4+41\cdot6+27\cdot13
    =1020=10\pmod{101}.
\]
Era possível a Alice fazer esta conta só com os seus segredos e o
desafio. Bob efectua a mesma combinação do lado dos pontos. Precisa
ainda das parcelas totalmente públicas, porque os pesos das matrizes
também entram nos vectores avaliados. Na primeira família, a
correcção é \(9(4,8)=(36,72)\). Na segunda, em bases
\(\boldsymbol h'\), é
\[
        (31,30)+9(2,8)+3(-6,-24)=(31,30).
\]
A escolha dos números fez aqui coincidir a última soma com a sua
primeira parcela. A expressão mostra de onde provieram as três,
para que uma coincidência aritmética do exemplo conserve o seu
devido lugar.

Podemos seguir cada coordenada sem escolher uma representação
concreta dos pontos da curva. Na família \(\boldsymbol g\), a
combinação que resulta dos pontos e da correcção é
\[
 9(3,4)+81(9,16)+41(1,2)+27(3,4)+(36,72)=(5,79).
\]
Para \(\boldsymbol h'\), lembremos que
\(h_1=h'_1\) e \(h_2=2h'_2\). Os coeficientes \((3,4)\) de
\(A_I\) nas bases originais passam a \((3,8)\) nas novas. Os
\((3,1)\) de \(S\) passam a \((3,2)\). As parcelas correspondentes
de \(A_O,K\) são nulas para a abertura tomada. Fica
\[
        9(3,8)+41(3,2)+(31,30)=(80,83).
\]
São os mesmos vectores que o provador encontrou ao avaliar
\(\boldsymbol l\) e \(\boldsymbol r\). As duas contas exprimem a
distributividade em direcções opostas. Alice avalia coordenadas
e só depois considera os respectivos geradores. Bob combina os
compromissos cujos geradores já transportam aquelas coordenadas.

Na igualdade do produto interno entram agora os pontos \(Q,L,R\).
A título de contabilidade, escrevamos que coeficiente de cada base
a abertura conhecida de Alice fornece a esses pontos. Cada linha
indica uma combinação formal das bases. Em particular,
\(U=\omega g\), com \(\omega\) dado pela transcrição desta fase.
O verificador calcula e soma pontos, sem precisar de decompor um
ponto recebido nem de encontrar qualquer logaritmo discreto.
\[
\begin{array}{c|rrrrr}
       &g_1&g_2&h'_1&h'_2&U\\ \hline
 Q     &5&79&80&83&89\\
 L     &0&5&83&0&11\\
 R     &79&0&0&80&58\\ \hline
 Q+4L+76R&50&99&8&2&97
\end{array}
\]
Por exemplo, \(L\) junta \(5g_2\), \(83h'_1\) e
\((5\cdot83)U=11U\). A terceira linha troca as metades,
produzindo \((79\cdot80)U=58U\). Com \(e=2\) temos \(e^2=4\)
e \(e^{-2}=76\), pelo que a linha final é exactamente a soma que
Bob pode formar na curva. As respostas \(a_*=100,b_*=4\) dão-lhe
a outra expressão da mesma linha. Os pesos de \(g_*\) são
\((51,2)\), e os de \(h_*\) são \((2,51)\). Logo,
\[
 100(51,2)=(50,99),\qquad 4(2,51)=(8,2),\qquad
 100\cdot4=97 .
\]
Todas as cinco colunas coincidiram. A última é a única em que
apareceu o produto das respostas. As quatro primeiras conservaram
a abertura dos vectores nas bases comprometidas.

Que aconteceria à descrição se uma das aberturas iniciais fosse
apresentada com outra máscara, tomando para ela o compromisso
correspondente? Na primeira abertura, substituir \(7\) por
\(7+\varepsilon\) na máscara de \(V\) desloca \(B_4\) por
\(32\varepsilon h\), e a máscara da avaliação por
\(81\cdot32\varepsilon=67\varepsilon\). O valor de \(v\) e os
cinco ligações lineares do circuito conservam-se. Na segunda
abertura, mudar só a máscara de \(K\) desloca \(P\) por
\(27\varepsilon h\) e \(\mu\) pelo escalar correspondente.
A subtracção de \(\mu h\) absorve o deslocamento. São identidades
para desafios fixados, que permitem ler a contribuição de cada
máscara. Numa execução de Fiat--Shamir, uma mudança nos compromissos
também muda os dados sobre os quais se calculam os desafios.

Assim se ligam as duas aberturas. Na igualdade do valor, \(89\)
aparece sobre \(g\), com a máscara \(83\) sobre \(h\). Na relação
vectorial que se vai dobrar, o mesmo \(89\) aparece sobre \(U\),
depois de retirada a máscara \(10\) sobre \(h\). A dobra conduz
ao produto final \(97\), e os termos cruzados explicam a diferença
entre os dois números. É possível seguir o valor por ambas as
contas mesmo conservando separados, em cada passo, o que Alice
conhece e os pontos que Bob efectivamente recebe.
